\SourceLeaf{187}and was now almost convinced that without some datum the Euclidean proposition could not be proved, and that he also thought it not unlikely that our geometry was only a chapter of a more general one. I then told him that some years ago you had publicly remarked that essentially no progress had been made in this matter since Euclid's time. Indeed, you had told me several times that, despite working extensively on the subject, you too had been unable to prove the absurdity of such a supposition. --- When he sent me the requested book, the accompanying slip was enclosed. Shortly afterward (at the end of December) he asked me in person to send you this slip when an opportunity arose and to ask you on his behalf to let him know your judgment of his ideas at your convenience.

Apart from everything else, the book itself has the agreeable feature of containing a rich bibliography of the subject, which, he tells me, he has continued to collect. .\ .\ .\ .\ .\}

\bigskip
\noindent\hrulefill
\bigskip

\Needspace{10\baselineskip}
\begin{center}[\emph{\textls{Enclosure:} Note by} \textsc{Schweikart}. \emph{Marburg, December 1818.}]\end{center}

\{There is a twofold geometry, --- a geometry in the narrower sense --- the \textls{Euclidean}, and an astral theory of magnitudes.

The triangles of the latter have the distinctive property that the sum of the three angles is not equal to two right angles.

Given \textls{this}, it can be proved with the utmost rigor that

a) the sum of the 3 angles in a triangle is \textls{less} than 2 right angles,

b) this sum becomes smaller the more area the triangle encloses,

c) the altitude of an isosceles right triangle continually increases the more its legs are extended, but it cannot exceed a certain line, which. I call the \textls{constant}.

